\chapter{How small a difference can we see: the paired object bootstrap}\label{ch:bootstrap}

\begin{onething}
The real PRISMA bench is 4 scenes and 909 objects, so a difference of 0.005 pAUC between two scorers can be luck in which
objects happen to be on it; the paired object bootstrap redraws the objects with replacement, the same draw for both
scorers, thousands of times, and since 2026-10-21 a gain counts only if the 95\% interval of that difference clears zero.
\end{onething}

\section{The problem: 909 objects is not many}\label{sec:bootstrap:problem}

By October 2026 every verdict on a new model or score leaned on gaps like 0.170 against 0.163 pAUC
(\cref{ch:measures} for pAUC). Nobody had asked how big a gap the bench can tell apart from noise. On 2026-10-21 that
question got its own small experiment and tool:

\begin{center}
\small
\begin{tabular}{p{5.6cm}p{7.0cm}}
\toprule
file & what it holds \\
\midrule
\file{research/experiments/bench-resolution/README.md} & the question, the method, the results, the new rule \\
\file{app/evaluation/significance.py} & \code{object_pauc_parts} and \code{paired_bootstrap}: the method in 63 lines \\
\file{scripts/compare_scorers.py} & the command line: two (model, scorer) pairs from a finished run \\
\file{tests/test_evaluation/test_significance.py} & two tests: parts reproduce pAUC; identical scorers differ by 0 \\
\bottomrule
\end{tabular}
\end{center}

The gains being chased were about 0.005 to 0.01 pAUC. FT04's real-bench pAUC with cloud excluded is
0.170\hfil\penalty50\hfilneg\tagM{research/experiments/bench-resolution/README.md}, so a gain of 0.005 is a change of
$0.005 / 0.170 = 0.029$, about 3\%. Would a bench of the same kind, with different objects on it, show the same 3\%?
You cannot collect another bench. You can ask the bench you have.

\begin{analogy}
Two teachers each mark the same class of 30 essays, and teacher B's average is 0.4 marks higher. Is B more generous, or
did a few essays happen to suit B? You cannot hire a new class, but you can ask ``if the class had been a slightly
different 30 essays, drawn from the same kind of students, would B still come out ahead?'' You redraw the class from the
30 you have, many times, and look at B minus A each time. Crucially, both teachers mark the \emph{same} redrawn class
each time; otherwise you would compare one teacher's easy class with the other's hard one.
\end{analogy}

\section{Sampling variation and the standard error}\label{sec:bootstrap:se}

Any number computed from data - a mean, a pAUC - would come out differently on different data of the same kind. The
spread of those hypothetical values is the number's \index{sampling distribution}\emph{sampling distribution}, and its
standard deviation is the \index{standard error}\emph{standard error}. For a plain mean there is a formula. You need two
rules, each shown on a tiny example first.

\paragraph{Rule 1: the variance of a sum.} For two random quantities $X$ and $Y$ with means $\mu_X$, $\mu_Y$, the
\index{covariance}\emph{covariance} is $\mathrm{Cov}(X,Y) = \E[(X - \mu_X)(Y - \mu_Y)]$: positive when they tend to
rise together. Then
\begin{align*}
\mathrm{Var}(X + Y) &= \E\big[(X + Y - \mu_X - \mu_Y)^2\big] && \text{definition of variance}\\
&= \E\big[((X - \mu_X) + (Y - \mu_Y))^2\big] && \text{regroup}\\
&= \E[(X - \mu_X)^2] + 2\,\E[(X - \mu_X)(Y - \mu_Y)] + \E[(Y - \mu_Y)^2] && \text{expand the square}\\
&= \mathrm{Var}\,X + 2\,\mathrm{Cov}(X, Y) + \mathrm{Var}\,Y.
\end{align*}
\emph{Tiny check}: let $X$ be a fair coin worth 0 or 1 and $Y = X$. Then $X + Y$ is 0 or 2 with equal chance, mean 1,
variance $\frac12(0-1)^2 + \frac12(2-1)^2 = 1$. The rule: $\mathrm{Var}\,X = \frac14$, $\mathrm{Cov}(X, X) = \frac14$,
so $\frac14 + 2 \times \frac14 + \frac14 = 1$. Agreed. If $X$ and $Y$ are \emph{independent}, the covariance is 0 and
the variances add.

\paragraph{Rule 2: the variance of a mean.} Take $n$ independent values, each of variance $\sigma^2$. Their mean is the
sum divided by $n$, and multiplying a quantity by $c$ multiplies its variance by $c^2$:
\begin{align*}
\mathrm{Var}\Big(\tfrac1n \textstyle\sum_i X_i\Big) &= \tfrac{1}{n^2}\,\mathrm{Var}\Big(\textstyle\sum_i X_i\Big)
  && \text{pull out } c = 1/n\\
&= \tfrac{1}{n^2} \cdot n\sigma^2 && \text{rule 1, independent: variances add}\\
&= \frac{\sigma^2}{n}, \qquad \text{standard error } = \frac{\sigma}{\sqrt n}.
\end{align*}

For a pAUC built from thousands of correlated pixels there is no such clean formula. The bootstrap gets the standard
error without one.

\section{The bootstrap idea, by hand}\label{sec:bootstrap:idea}

\begin{definition}[Bootstrap]
A \index{bootstrap}\emph{bootstrap} \citep{efron1979bootstrap} estimates the sampling distribution of a number using
only the data you have. From the $n$ data items, draw $n$ items \emph{with replacement} - each draw picks any item, so
some appear twice and some not at all - recompute the number on that \emph{bootstrap sample}, and repeat many times. The
spread of the recomputed numbers stands in for the spread you would have seen across fresh data.
\end{definition}

The logic: the data are your best picture of the population they came from. Drawing from the data with replacement is
drawing from that picture, so it mimics collecting a new dataset.

\begin{worked}[label={we:bootstrap:five}]{Bootstrapping the mean of five numbers}
Do this on paper. The data are $x = (2, 4, 5, 7, 12)$, with mean $(2 + 4 + 5 + 7 + 12)/5 = 30/5 = 6$. Six bootstrap
draws are given\tagA{} as positions 1--5 (a random number generator would supply them):
\begin{center}
\begin{tabular}{lccc}
\toprule
draw & positions & values & mean \\
\midrule
1 & 4, 1, 5, 5, 2 & 7, 2, 12, 12, 4 & $37/5 = 7.4$ \\
2 & 2, 2, 3, 1, 4 & 4, 4, 5, 2, 7 & $22/5 = 4.4$ \\
3 & 3, 5, 1, 3, 2 & 5, 12, 2, 5, 4 & $28/5 = 5.6$ \\
4 & 1, 1, 2, 4, 3 & 2, 2, 4, 7, 5 & $20/5 = 4.0$ \\
5 & 5, 4, 3, 2, 5 & 12, 7, 5, 4, 12 & $40/5 = 8.0$ \\
6 & 2, 3, 4, 4, 1 & 4, 5, 7, 7, 2 & $25/5 = 5.0$ \\
\bottomrule
\end{tabular}
\end{center}
Draw 1 contains 12 twice and 5 not at all: that is ``with replacement''. The mean of the six bootstrap means:
\begin{align*}
(7.4 + 4.4 + 5.6 + 4.0 + 8.0 + 5.0)/6 = 34.4/6 = 5.733.
\end{align*}
Their spread (population standard deviation): deviations from 5.733 are $1.667, -1.333, -0.133, -1.733, 2.267, -0.733$;
squares $2.778, 1.778, 0.018, 3.004, 5.138, 0.538$; sum $13.253$; divided by 6, $2.209$; square root
\[
\text{bootstrap standard error} \approx 1.486.
\]
Compare rule 2: the data's population standard deviation is $\sqrt{(16 + 4 + 1 + 1 + 36)/5} = \sqrt{58/5} =
\sqrt{11.6} = 3.406$, so $\sigma/\sqrt5 = 3.406/2.236 = 1.523$. Six draws already land near it; with 100,000 draws the
book's script gets 1.522\tagC{}.
\end{worked}

\section{The percentile interval}\label{sec:bootstrap:percentile}

\begin{definition}[Percentile 95\% interval]
The \index{percentile interval}\emph{percentile 95\% interval} from a bootstrap is the 2.5th and 97.5th percentiles of
the recomputed numbers: the range holding the middle 95\% of them.
\end{definition}

\code{paired_bootstrap} uses \code{np.percentile(..., [2.5, 97.5])}, which interpolates linearly. With $n$ values sorted
from smallest (position 0) to largest (position $n - 1$), the $q$-th percentile sits at position $(n - 1)\,q/100$; if
that falls between two positions, it takes the matching share of the gap.

\begin{worked}[label={we:bootstrap:pct}]{Percentiles of the six bootstrap means}
Sorted: $4.0, 4.4, 5.0, 5.6, 7.4, 8.0$ at positions 0--5.
\begin{align*}
\text{2.5th: position } (6 - 1) \times 0.025 &= 0.125 && \text{between positions 0 and 1}\\
4.0 + 0.125 \times (4.4 - 4.0) &= 4.0 + 0.05 = 4.05\\
\text{97.5th: position } 5 \times 0.975 &= 4.875 && \text{between positions 4 and 5}\\
7.4 + 0.875 \times (8.0 - 7.4) &= 7.4 + 0.525 = 7.925.
\end{align*}
The interval from six draws is $[4.05, 7.925]$; \code{np.percentile} returns the same\tagC{}. Six draws are far too
few to trust; 100,000 give $[3.4, 9.2]$\tagC{}. With 1,000 draws the 2.5th percentile sits at position
$999 \times 0.025 = 24.975$ (between the 25th and 26th smallest) and the 97.5th at $974.025$; with 2,000 at $49.975$
and $1{,}949.025$.
\end{worked}

\paragraph{The whole bootstrap distribution, exactly.} Five items drawn five times with replacement give
$5^5 = 5 \times 5 \times 5 \times 5 \times 5 = 3{,}125$ equally likely ordered draws, few enough to list them all. The
book's script does, and \cref{fig:bootstrap:exact} plots how often each bootstrap mean occurs. Its standard deviation is
1.5232, exactly $\sigma/\sqrt5$, and its 2.5th and 97.5th percentiles are 3.4 and 9.2\tagC{}: the 100,000 random draws
of \cref{we:bootstrap:pct} had already found the exact answer. Real benches have far too many objects to enumerate, which
is why the code draws at random, 1,000 or 2,000 times.

\begin{figure}[htbp]
\centering
\begin{tikzpicture}
\begin{axis}[width=12cm,height=5.6cm,ybar,bar width=4pt,xmin=1.6,xmax=12.4,ymin=0,ymax=230,
  xlabel={bootstrap mean of $(2, 4, 5, 7, 12)$},ylabel={draws out of 3,125},
  xtick={2,3,4,5,6,7,8,9,10,11,12},tick label style={font=\small}]
\fill[spec!12] (axis cs:3.4,0) rectangle (axis cs:9.2,230);
\addplot[fill=spat!60,draw=spat] coordinates {(2,1) (2.4,5) (2.6,5) (2.8,10) (3,25) (3.2,20) (3.4,50) (3.6,55)
 (3.8,60) (4,106) (4.2,75) (4.4,125) (4.6,125) (4.8,105) (5,181) (5.2,110) (5.4,175) (5.6,175) (5.8,120) (6,215)
 (6.2,110) (6.4,160) (6.6,155) (6.8,80) (7,171) (7.2,70) (7.4,110) (7.6,110) (7.8,40) (8,105) (8.2,30) (8.4,50)
 (8.6,50) (8.8,10) (9,50) (9.2,10) (9.4,20) (9.6,20) (10,15) (10.4,5) (10.6,5) (11,5) (12,1)};
\draw[loss,thick,dashed] (axis cs:6,0) -- (axis cs:6,230);
\node[font=\small,text=loss,anchor=west] at (axis cs:6.05,218) {data mean 6};
\node[font=\small,text=spec!70!black,anchor=west,align=left] at (axis cs:9.9,150) {95\% interval\\$[3.4, 9.2]$};
\end{axis}
\end{tikzpicture}
\caption{The exact bootstrap distribution of the mean of five numbers: every one of the $5^5 = 3{,}125$ draws with
replacement, counted by the mean it gives (computed by \texttt{research/textbook/tools/compute\_bootstrap.py}). The teal
band is the percentile 95\% interval. The distribution is lumpy because only 43 distinct means are possible, and skewed
right because one value (12) is far from the rest.}
\label{fig:bootstrap:exact}
\end{figure}

\section{Why resample objects, not pixels}\label{sec:bootstrap:objects}

pAUC averages over target pixels, but target pixels are not independent pieces of evidence. The pixels of one object rise
and fall together: if the model reconstructs a rooftop well, all its pixels score low at once. Swap a dozen hard objects
for a dozen easy ones and pAUC moves; swap a dozen pixels inside one object and it barely does.

\paragraph{How much dependence costs, derived.} Model $n$ objects of $m$ pixels each. Each pixel's value (think of its
contribution $c$ below) has variance $\sigma^2$; two pixels of the same object have correlation $\rho$, so covariance
$\rho\sigma^2$; pixels of different objects are independent. Apply rule 1 to one object's $m$ pixels: $m$ variance
terms and $m(m-1)$ ordered pairs of distinct pixels, each a covariance:
\begin{align*}
\mathrm{Var}(\text{object sum}) &= m\sigma^2 + m(m-1)\rho\sigma^2 = m\sigma^2\,[1 + (m-1)\rho].\\
\intertext{Objects are independent, so their sums' variances add, and the mean over all $nm$ pixels divides by $nm$:}
\mathrm{Var}(\text{pixel mean}) &= \frac{n \cdot m\sigma^2\,[1 + (m-1)\rho]}{(nm)^2}
  = \frac{\sigma^2}{nm}\,[1 + (m-1)\rho].
\end{align*}
Compare rule 2 for $N$ independent pixels, $\sigma^2/N$. The two agree when
\begin{equation}\label{eq:bootstrap:neff}
N_{\text{eff}} = \frac{nm}{1 + (m-1)\rho},
\end{equation}
the \index{effective sample size}\emph{effective sample size}; the bracket $1 + (m-1)\rho$ is the \emph{design effect}.
At $\rho = 0$ every pixel counts; at $\rho = 1$, $N_{\text{eff}} = nm/m = n$: you have as much evidence as you have
\emph{objects}.

\begin{worked}[label={we:bootstrap:neff}]{An effective sample size}
Suppose\tagA{} 909 objects of 4 pixels each ($909 \times 4 = 3{,}636$ target pixels) with within-object correlation
$\rho = 0.8$.
\begin{align*}
\text{design effect} &= 1 + (4 - 1) \times 0.8 = 1 + 2.4 = 3.4\\
N_{\text{eff}} &= 3{,}636 / 3.4 = 1{,}069.4\\
\text{standard error ratio} &= \sqrt{3.4} = 1.84.
\end{align*}
3,636 pixels carry the evidence of about 1,069 independent ones. A pixel bootstrap, which treats them as independent,
would report a standard error $1.84$ times too small.
\end{worked}

The book's script checks this by simulation\tagA{}: 300 objects of sizes 1--10 (967 pixels), each object's pixels
sharing one value plus small noise. Resampling objects gave a standard error of 0.0235, resampling pixels 0.0098: the
pixel bootstrap was 2.4 times too confident\tagC{}. That is the reason in \file{significance.py}'s module docstring for
drawing ``each scene's objects with replacement''.

\begin{keypoint}
More numbers is not more evidence: resample the unit that varies independently, here the object.
\end{keypoint}

\paragraph{The background is held fixed.} Each scene's background pixels are not resampled. The docstring's reason:
``there are millions; their sampling error is small next to the objects'''. Only the target side of the measure is
redrawn.

\section{pAUC rebuilt from per-pixel contributions}\label{sec:bootstrap:cp}

To redraw objects you need pAUC as a sum over objects. The docstring of \file{app/evaluation/significance.py} gives the
rearrangement.

\begin{definition}[Per-pixel contribution]
For a target pixel $p$, let $\mathrm{FPR}(p)$ be the share of the scene's background pixels scoring \emph{at or above}
it. With limit $m$ (\code{max_fpr}, default 0.01), its \index{per-pixel contribution}\emph{contribution} is
\begin{equation}\label{eq:bootstrap:c}
c(p) = \frac{\max(0,\ m - \mathrm{FPR}(p))}{m},
\end{equation}
and the scene's pAUC is the mean of $c$ over its $P$ target pixels:
$\mathrm{pAUC}_m = \frac1P\sum_p c(p)$.
\end{definition}

Read $c(p)$ as ``the share of the first $m$ of false alarms for which this pixel is caught''. A target pixel above every
background pixel has $\mathrm{FPR} = 0$ and $c = 1$; one that beats the bottom 99.5\% has $\mathrm{FPR} = 0.005$ and is
caught over the second half of the range, $c = 0.5$; one at or below the top 1\% of background has $c = 0$.

\paragraph{Proof that the mean of $c$ is the pAUC (no ties).} In the ROC walk of \cref{ch:measures}, target $p$ is
passed when the false-alarm rate reaches $\mathrm{FPR}(p)$, and from then on it counts in the TPR. So at false-alarm
rate $f$,
\[
\mathrm{TPR}(f) = \frac1P \sum_p \mathbf{1}[\mathrm{FPR}(p) \le f].
\]
Put this into the definition of pAUC (\cref{eq:measures:pauc}):
\begin{align*}
\mathrm{pAUC}_m &= \frac1m \int_0^m \frac1P \sum_p \mathbf{1}[\mathrm{FPR}(p) \le f]\,df && \text{substitute TPR}\\
&= \frac1P \sum_p \frac1m \int_0^m \mathbf{1}[\mathrm{FPR}(p) \le f]\,df && \text{swap the sum and the integral}\\
&= \frac1P \sum_p \frac{\max(0,\ m - \mathrm{FPR}(p))}{m} && \text{the indicator is 1 on } [\mathrm{FPR}(p), m]\\
&= \frac1P \sum_p c(p).
\end{align*}
The third line: the indicator is 1 for $f$ from $\mathrm{FPR}(p)$ up to $m$, an interval of length
$m - \mathrm{FPR}(p)$, or of length 0 if $\mathrm{FPR}(p) \ge m$. With tied scores \code{exact_roc} draws a diagonal
segment while $c(p)$ treats every pixel as its own step, so the two agree ``up to ties''; with millions of distinct float
scores the gap is negligible, and the README reports the bootstrap's point values reproducing the evaluator's (0.1621 and
0.1695).

So an object needs only two numbers: the \emph{sum of $c$ over its pixels} and its \emph{pixel count}.
\code{object_pauc_parts} returns them:
\begin{lstlisting}[basicstyle=\ttfamily\footnotesize]
bg = np.sort(score[valid & (objects == 0)])                      # (N,) bg scores, ascending
t = valid & (objects > 0)                                         # (H, W) target pixels
fpr = 1.0 - np.searchsorted(bg, score[t], side="left") / max(len(bg), 1)   # share of bg >= score
c = np.maximum(0.0, max_fpr - fpr) / max_fpr                      # (P,) eq. c
ids = objects[t]
sums, counts = np.bincount(ids, weights=c), np.bincount(ids)      # per object: sum of c, count
\end{lstlisting}
\code{searchsorted(..., side="left")} counts the background scores strictly below each target score, so one minus that
share is the share at or above it. A bootstrap draw adds up the drawn objects' pairs and divides; no score map is touched
again, which is why thousands of draws are cheap.

\begin{worked}[label={we:bootstrap:tiny}]{A tiny scene, both ways}
One scene, 1,000 background pixels, three objects, four target pixels\tagA{}:
\begin{align*}
\text{obj 1, pixel a: } & 2 \text{ bg above}, & \mathrm{FPR} &= 0.002, & c &= (0.01 - 0.002)/0.01 = 0.8\\
\text{obj 1, pixel b: } & 6 \text{ bg above}, & \mathrm{FPR} &= 0.006, & c &= (0.01 - 0.006)/0.01 = 0.4\\
\text{obj 2, pixel c: } & 0 \text{ bg above}, & \mathrm{FPR} &= 0.000, & c &= 0.01/0.01 = 1.0\\
\text{obj 3, pixel d: } & 30 \text{ bg above}, & \mathrm{FPR} &= 0.030, & c &= \max(0, -0.02)/0.01 = 0
\end{align*}
Object parts: object 1 $(0.8 + 0.4, 2) = (1.2, 2)$; object 2 $(1.0, 1)$; object 3 $(0.0, 1)$.
\[
\mathrm{pAUC} = \frac{1.2 + 1.0 + 0.0}{2 + 1 + 1} = \frac{2.2}{4} = 0.55.
\]
Now the ROC area. Sort the target pixels by FPR: TPR is $\frac14$ from FPR 0, $\frac24$ from 0.002, $\frac34$ from
0.006, and pixel d never arrives before 0.01:
\begin{align*}
0 \text{ to } 0.002: &\quad 0.002 \times 0.25 = 0.0005\\
0.002 \text{ to } 0.006: &\quad 0.004 \times 0.50 = 0.0020\\
0.006 \text{ to } 0.010: &\quad 0.004 \times 0.75 = 0.0030\\
\text{area} &= 0.0005 + 0.0020 + 0.0030 = 0.0055, \qquad 0.0055 / 0.01 = 0.55.
\end{align*}
The same 0.55. The book's script builds this scene and gets 0.55 from \code{object_pauc_parts} and from
\code{exact_roc} plus \code{pauc}\tagC{}.
\end{worked}

\section{The paired design}\label{sec:bootstrap:paired}

To compare scorers A and B, you want the interval of the \emph{difference} $\mathrm{pAUC}_B - \mathrm{pAUC}_A$. You could
bootstrap A and B separately, each with its own draws. That throws away the most useful fact: both scorers are judged on
the same objects, and an object that is hard for A is usually hard for B.

\paragraph{The variance of a difference, derived.} Apply rule 1 with $Y$ replaced by $-B$ (and
$\mathrm{Var}(-B) = \mathrm{Var}\,B$, $\mathrm{Cov}(A, -B) = -\mathrm{Cov}(A, B)$), or directly:
\begin{align*}
\mathrm{Var}(A - B) &= \E\big[((A - B) - (\mu_A - \mu_B))^2\big] && \text{definition}\\
&= \E\big[((A - \mu_A) - (B - \mu_B))^2\big] && \text{regroup}\\
&= \E[(A - \mu_A)^2] - 2\,\E[(A - \mu_A)(B - \mu_B)] + \E[(B - \mu_B)^2] && \text{expand the square}\\
&= \mathrm{Var}\,A + \mathrm{Var}\,B - 2\,\mathrm{Cov}(A, B).
\end{align*}
When A and B rise and fall together, the covariance is large and positive, and it cancels most of the variance. Drawing
independently sets that covariance to zero by construction, and you get $\mathrm{Var}\,A + \mathrm{Var}\,B$.

\begin{worked}[label={we:bootstrap:var}]{Five objects, paired and unpaired}
Per-object scores\tagA{} for two scorers on the same five objects:
$A = (0.9, 0.2, 0.6, 0.1, 0.7)$, $B = (0.95, 0.25, 0.6, 0.15, 0.75)$.
\begin{align*}
\mu_A &= (0.9 + 0.2 + 0.6 + 0.1 + 0.7)/5 = 2.5/5 = 0.50\\
A - \mu_A &= (0.40, -0.30, 0.10, -0.40, 0.20)\\
\mathrm{Var}\,A &= (0.16 + 0.09 + 0.01 + 0.16 + 0.04)/5 = 0.46/5 = 0.0920\\
\mu_B &= (0.95 + 0.25 + 0.6 + 0.15 + 0.75)/5 = 2.7/5 = 0.54\\
B - \mu_B &= (0.41, -0.29, 0.06, -0.39, 0.21)\\
\mathrm{Var}\,B &= (0.1681 + 0.0841 + 0.0036 + 0.1521 + 0.0441)/5 = 0.452/5 = 0.0904\\
\mathrm{Cov}(A,B) &= (0.164 + 0.087 + 0.006 + 0.156 + 0.042)/5 = 0.455/5 = 0.0910\\
\mathrm{Var}(A - B) &= 0.0920 + 0.0904 - 2 \times 0.0910 = 0.1824 - 0.1820 = 0.0004.
\end{align*}
Check directly: $A - B = (-0.05, -0.05, 0, -0.05, -0.05)$, mean $-0.04$, deviations
$(-0.01, -0.01, 0.04, -0.01, -0.01)$, variance $(4 \times 0.0001 + 0.0016)/5 = 0.002/5 = 0.0004$. Agreed. Unpaired, the
variance would be $0.0920 + 0.0904 = 0.1824$, $456$ times larger\tagC{}: B's consistent $+0.05$ edge would vanish in the
noise of which objects were drawn.
\end{worked}

\code{paired_bootstrap} therefore draws one set of object positions per scene and applies it to both scorers.

\begin{worked}[label={we:bootstrap:draws}]{Two paired draws on the tiny scene}
Give scorer B its own parts on the same objects of \cref{we:bootstrap:tiny}: object 1 $(1.0, 2)$, object 2 $(0.9, 1)$,
object 3 $(0.2, 1)$. Point values: $A = 2.2/4 = 0.550$, $B = (1.0 + 0.9 + 0.2)/4 = 2.1/4 = 0.525$, $B - A = -0.025$.

Draw objects 1, 1, 3 (with replacement), the same for both:
\begin{align*}
A &= (1.2 + 1.2 + 0.0)/(2 + 2 + 1) = 2.4/5 = 0.48\\
B &= (1.0 + 1.0 + 0.2)/(2 + 2 + 1) = 2.2/5 = 0.44, \qquad B - A = -0.04.
\end{align*}
Draw objects 2, 2, 2:
\begin{align*}
A &= (1.0 + 1.0 + 1.0)/3 = 1.00, \qquad B = (0.9 + 0.9 + 0.9)/3 = 0.90, \qquad B - A = -0.10.
\end{align*}
A swung from 0.48 to 1.00 between the two draws; $B - A$ moved only from $-0.04$ to $-0.10$. The draw moves both scorers
together, and the difference cancels most of it.

With only three objects the interval is still wide: \code{paired_bootstrap} with 1,000 draws on this toy gives
$[-0.10, +0.20]$ and $P(B > A) = 0.28$\tagC{}. (A draw of object 3 three times gives $B - A = 0.2 - 0 = +0.2$.) Few
objects, wide interval: that is the whole problem in miniature.
\end{worked}

\paragraph{Scenes and the two measures.} On the real bench each scene's objects are drawn separately, with
\code{rng.integers(len(ids[s]), size=len(ids[s]))} per scene; each scene's pAUC is computed from its drawn parts, and
the four are averaged - the same ``mean over scenes'' as the evaluator. If per-object hits are given (1 if caught at a
cut-off, from the run's \file{per_object.csv}, column \code{hit_0.1} by default), objects caught is resampled too, but
\emph{pooled} across scenes: all caught drawn objects over all drawn objects. The function also returns
$P(B > A)$, the share of draws in which B's pAUC beat A's.

\section{What it does not cover: only four scenes}\label{sec:bootstrap:scenes}

The bootstrap redraws objects \emph{inside} each scene and never redraws the scenes. A model can be lucky in the scenes
too. Why not bootstrap scenes as well? Count the distinct ways to draw 4 scenes from 4 with replacement, order not
mattering. That is choosing 4 items from 4 kinds with repetition, $\binom{4 + 4 - 1}{4}$:
\begin{align*}
\binom{7}{4} &= \frac{7 \times 6 \times 5 \times 4}{4 \times 3 \times 2 \times 1} = \frac{840}{24} = 35.
\end{align*}
Thirty-five possible draws, one of which is ``scene 2 four times'': too few and too lumpy for an interval worth trusting.
The README says it plainly: ``Scene-level uncertainty (4 scenes) is not covered and is larger: this interval is a lower
bound on the real uncertainty.'' Nor does it cover checkpoint-to-checkpoint variation: FT09c, five more epochs of the
same run as FT09b, moved by as much as FT09b's resolved gain (\cref{tab:bootstrap:results}).

\section{The rule, and the intervals so far}\label{sec:bootstrap:rule}

\begin{keypoint}
A gain counts only if its paired-bootstrap 95\% interval clears zero (rule adopted 2026-10-21).
\end{keypoint}

Run it on a finished evaluation:
\begin{lstlisting}[basicstyle=\ttfamily\footnotesize,language={}]
python scripts/compare_scorers.py configs/eval/scene_adapt.json "FT04_ep10|RM, cloud excluded" \
    "adapt50|RM, cloud excluded" [--top 0.1] [--boot 2000]
\end{lstlisting}
Each side is \texttt{<model name>|<scorer label>} as in the run's \file{per_scene.csv} (\texttt{-} as the model name for
a model-free scorer). It reads the cached score maps and \file{per_object.csv} of the \emph{finished} run, only for
\code{prisma_bench} scenes, recomputes nothing, and prints both point values, $B - A$, the 95\% interval and $P(B>A)$.

\begin{sourcenote}[: how many draws]
The README states its results used 1,000 draws. The code's defaults are 2,000: \code{paired_bootstrap(..., n=2000,
seed=0)} in \file{app/evaluation/significance.py} (line 32) and \texttt{-{}-boot 2000} in
\file{scripts/compare_scorers.py} (line 37). So the README's runs passed \texttt{-{}-boot 1000}; re-running with the
defaults will give slightly different interval ends. The generator is seeded (\code{seed=0}), so a given $n$ is
reproducible.
\end{sourcenote}

\begin{drift}
Paired object bootstrap on the real PRISMA bench, pAUC $\le$ 1\%, $B - A$ [95\% interval]. Rows 1--3 from
\file{research/experiments/bench-resolution/README.md} (2026-10-21, 1,000 draws); rows 4--6 from the run table of
\file{docs/lld/evaluation.md} (2026-10-22/23). FT04c is FT04 with cloud excluded; gated adaptation is $N = 50$.
\begin{center}
\small
\begin{tabular}{clcc}
\toprule
 & A $\to$ B & pAUC $B - A$ [95\%] & caught 0.1\%, pts [95\%] \\
\midrule
1 & FT04 $\to$ + cloud excl. & $+0.0074$ [$+0.0047$, $+0.0101$] & $+1.0$ [$+0.4$, $+1.7$] \\
2 & FT04c $\to$ gated adapt. & $-0.0091$ [$-0.0176$, $-0.0028$] & $-0.7$ [$-1.7$, $+0.2$] \\
3 & FT04c $\to$ FT06 $/\sqrt{b}$ & $-0.0067$ [$-0.0174$, $+0.0037$] & $+0.8$ [$-0.4$, $+2.0$] \\
4 & FT04c $\to$ FT09b & $+0.0060$ [$+0.0002$, $+0.0118$] & unresolved \\
5 & FT04c $\to$ FT09c & $+0.0020$ [$-0.0026$, $+0.0064$] & - \\
6 & FT09b $\to$ FT09c & $-0.0040$ [$-0.0076$, $-0.0010$] & - \\
\bottomrule
\end{tabular}
\captionof{table}{Paired-bootstrap intervals on the real PRISMA bench.}\label{tab:bootstrap:results}
\end{center}
\end{drift}

\begin{figure}[htbp]
\centering
\begin{tikzpicture}
\begin{axis}[width=12cm,height=7.2cm,xmin=-0.02,xmax=0.015,ymin=0.4,ymax=6.6,
  xlabel={pAUC difference $B - A$ (95\% interval)},
  ytick={1,2,3,4,5,6},
  yticklabels={{FT09b $\rightarrow$ FT09c},{FT04 $\rightarrow$ FT09c},{FT04 $\rightarrow$ FT09b},{FT04 $\rightarrow$ FT06 $\sqrt{b}$},{FT04 $\rightarrow$ gated adapt.},{FT04 $\rightarrow$ + cloud excl.}},
  yticklabel style={font=\small},xticklabel style={font=\small,/pgf/number format/fixed,/pgf/number format/precision=3},
  scaled x ticks=false,xtick={-0.02,-0.015,-0.01,-0.005,0,0.005,0.01,0.015},grid=major,grid style={gray!15}]
\draw[black,thick] (axis cs:0,0.4) -- (axis cs:0,6.6);
% row 6: cloud exclusion, resolved gain
\draw[spec,line width=2pt] (axis cs:0.0047,6) -- (axis cs:0.0101,6);
\addplot[only marks,mark=*,spec,mark size=2.5pt] coordinates {(0.0074,6)};
% row 5: gated adaptation, resolved loss
\draw[loss,line width=2pt] (axis cs:-0.0176,5) -- (axis cs:-0.0028,5);
\addplot[only marks,mark=*,loss,mark size=2.5pt] coordinates {(-0.0091,5)};
% row 4: FT06 sqrt b, tie
\draw[maskgrey,line width=2pt] (axis cs:-0.0174,4) -- (axis cs:0.0037,4);
\addplot[only marks,mark=*,maskgrey,mark size=2.5pt] coordinates {(-0.0067,4)};
% row 3: FT09b, resolved gain
\draw[spec,line width=2pt] (axis cs:0.0002,3) -- (axis cs:0.0118,3);
\addplot[only marks,mark=*,spec,mark size=2.5pt] coordinates {(0.0060,3)};
% row 2: FT09c vs FT04, tie
\draw[maskgrey,line width=2pt] (axis cs:-0.0026,2) -- (axis cs:0.0064,2);
\addplot[only marks,mark=*,maskgrey,mark size=2.5pt] coordinates {(0.0020,2)};
% row 1: FT09c vs FT09b, resolved loss
\draw[loss,line width=2pt] (axis cs:-0.0076,1) -- (axis cs:-0.0010,1);
\addplot[only marks,mark=*,loss,mark size=2.5pt] coordinates {(-0.0040,1)};
\end{axis}
\end{tikzpicture}
\caption{The paired-bootstrap intervals of \cref{tab:bootstrap:results}, drawn to scale: point value (dot) and 95\%
interval (bar) of the pAUC difference $B - A$. Teal: interval wholly above zero (a resolved gain); amber: wholly below
zero (a resolved loss); grey: crosses zero (a tie). FT09b's interval clears zero by 0.0002.}
\label{fig:bootstrap:intervals}
\end{figure}

\paragraph{Reading the rows.}
\begin{itemize}
\item \textbf{Cloud exclusion is a real gain} on both measures (row 1): the one change adopted since FT04 at the time.
Its interval is the narrowest, half-width $(0.0101 - 0.0047)/2 = 0.0027$. A likely reason, from the method rather than
the README: excluding cloud leaves most objects' scores untouched, so most per-object differences are exactly zero and the
paired difference barely wobbles - \cref{we:bootstrap:var}'s covariance at its largest.
\item \textbf{Gated per-scene adaptation's loss is real} on pAUC (row 2), though its objects-caught interval straddles
zero. Read each measure's interval on its own.
\item \textbf{FT06's square-root variant is a tie} (row 3): both of its intervals cross zero, whatever the point values say.
\item \textbf{FT09b was the first resolved gain since cloud exclusion} (row 4), by the narrowest of margins; and FT09c,
five epochs later, is a tie with FT04 and a resolved loss against FT09b (rows 5, 6). The object bootstrap does not cover
that checkpoint-to-checkpoint variation.
\end{itemize}

The README summarises the bench's resolution as interval half-widths of about $\pm 0.005$ to $\pm 0.01$ pAUC and $\pm 1$ to
$\pm 1.5$ points of objects caught at the top 0.1\%. So an old pass criterion of ``+0.005 pAUC'' sits at the edge of what
this bench can resolve: a run that clears it on its point value may still be a tie.

\begin{confusions}
\confusion{The interval says there is a 95\% chance the true gain lies inside it?}{It says how far the difference wobbles
when the \emph{objects} are redrawn, scenes and checkpoints held fixed. Scene luck is extra, so the real uncertainty is
wider: treat the interval as a lower bound.}
\confusion{Resampling pixels would be more precise, since there are far more of them?}{More numbers, not more evidence.
An object's pixels move together; counting them separately manufactures confidence (a design effect of 3.4 makes the
standard error 1.84 times too small).}
\confusion{Why not bootstrap A and B separately and see if their intervals overlap?}{Because the draw moves both together.
Separate intervals would overlap almost always, and you would call real differences ties
(\cref{we:bootstrap:var}: variance 0.1824 unpaired against 0.0004 paired).}
\confusion{pAUC from $c(p)$ is an approximation of the real thing?}{It is the same area, rearranged so each target pixel's
share can be added up. It differs from the evaluator only at tied scores.}
\confusion{Objects caught is averaged per scene, like pAUC?}{Not in the bootstrap: pAUC is a mean over the four scenes;
objects caught is pooled over all drawn objects.}
\confusion{The gated adaptation loss straddles zero, so it is a tie?}{On objects caught, yes. On pAUC its interval
$[-0.0176, -0.0028]$ is wholly below zero, so the pAUC loss is real.}
\end{confusions}

\begin{checkyourself}
\item\label{ex:bootstrap:1} In \cref{we:bootstrap:tiny}, pixel d's score rises until 5 background pixels score at or above
it. What is its $c$, and the scene's pAUC for A?
\item\label{ex:bootstrap:2} Bootstrap the mean of $(1, 3, 8)$ by hand with the four given draws (positions)
$(1,1,2)$, $(3,2,3)$, $(2,2,1)$, $(3,3,3)$. Give the four means, their mean, and the percentile interval
\code{np.percentile} would return for $[2.5, 97.5]$.
\item\label{ex:bootstrap:3} A result reads $+0.006$ [$-0.001$, $+0.013$]. Does it count as a gain under the rule? What is
its half-width?
\item\label{ex:bootstrap:4} Two scorers on four objects: $A = (0.8, 0.3, 0.5, 0.4)$, $B = (0.9, 0.3, 0.6, 0.6)$. Compute
$\mathrm{Var}\,A$, $\mathrm{Var}\,B$, $\mathrm{Cov}(A,B)$ and $\mathrm{Var}(A - B)$ by the formula, and check
$\mathrm{Var}(A - B)$ directly.
\item\label{ex:bootstrap:5} 909 objects average 2 pixels each with within-object correlation 0.9. What are the design
effect and the effective sample size? By what factor would a pixel bootstrap understate the standard error?
\item\label{ex:bootstrap:6} In the paired example, draw objects 1, 2, 3 (each once). Compute A, B and $B - A$. Why is
this the point value?
\item\label{ex:bootstrap:7} With 2,000 bootstrap draws, at which 0-based positions do the 2.5th and 97.5th percentiles
sit, and between which sorted draws (counting from 1) do they interpolate?
\end{checkyourself}

\nextchapter{ch:sortbug}
